% ══════════════════════════════════════════════════════════════
%  The OMARWA–E₈ Correspondence: A Research Program
%  Connecting Modular Periodicity to Exceptional Lie Structure
%
%  Paket 4 — E₈ Research Program Preprint (v1.0)
%  Author: Ömer Çetintaş
%  Date:   March 2026
% ══════════════════════════════════════════════════════════════
\documentclass[11pt, a4paper]{article}

% ── Encoding & Fonts ──────────────────────────────────────────
\usepackage[utf8]{inputenc}
\usepackage[T1]{fontenc}
\usepackage{lmodern}
\usepackage{microtype}

% ── Page Geometry ─────────────────────────────────────────────
\usepackage[
  a4paper,
  top=2.5cm,
  bottom=2.5cm,
  left=2.5cm,
  right=2.5cm,
  headheight=14pt
]{geometry}

% ── Mathematics ───────────────────────────────────────────────
\usepackage{amsmath,amssymb,amsthm,mathtools}

% ── Colors ────────────────────────────────────────────────────
\usepackage[dvipsnames]{xcolor}

% ── Tables & Floats ──────────────────────────────────────────
\usepackage{booktabs}
\usepackage{array}
\usepackage{float}
\usepackage[font=small, labelfont=bf, labelsep=period]{caption}

% ── Header / Footer ──────────────────────────────────────────
\usepackage{fancyhdr}
\pagestyle{fancy}
\fancyhf{}
\fancyhead[L]{\small\textit{The OMARWA--$E_8$ Correspondence}}
\fancyhead[R]{\small\textit{\c{C}etinta\c{s}, 2026}}
\fancyfoot[C]{\small\thepage}
\renewcommand{\headrulewidth}{0.4pt}
\renewcommand{\footrulewidth}{0pt}

\fancypagestyle{firstpage}{%
  \fancyhf{}
  \fancyfoot[C]{\small\thepage}
  \renewcommand{\headrulewidth}{0pt}
}

% ── Hyperlinks ────────────────────────────────────────────────
\usepackage[
  colorlinks=true,
  linkcolor=blue!70!black,
  citecolor=blue!70!black,
  urlcolor=blue!70!black,
  pdfauthor={\"Omer \c{C}etinta\c{s}},
  pdftitle={The OMARWA-E8 Correspondence: A Research Program},
  pdfsubject={Mathematical Physics, Lie Algebras, Number Theory},
  pdfkeywords={E8, OMARWA sequence, Coxeter number, Weyl exponents, P622, research program}
]{hyperref}

% ── Theorem Environments ─────────────────────────────────────
\theoremstyle{plain}
\newtheorem{theorem}{Theorem}[section]
\newtheorem{lemma}[theorem]{Lemma}
\newtheorem{proposition}[theorem]{Proposition}
\newtheorem{corollary}[theorem]{Corollary}

\theoremstyle{definition}
\newtheorem{definition}[theorem]{Definition}
\newtheorem{observation}{Observation}[section]

\theoremstyle{remark}
\newtheorem{remark}{Remark}[section]

% ── Custom Commands ───────────────────────────────────────────
\newcommand{\ord}{\operatorname{ord}}
\DeclareMathOperator{\lcm}{lcm}
\DeclareMathOperator{\DR}{DR}
\DeclareMathOperator{\rank}{rank}

% ── Spacing ───────────────────────────────────────────────────
\usepackage{enumitem}
\setlist[enumerate]{leftmargin=2em, itemsep=0.3em}
\setlist[itemize]{leftmargin=2em, itemsep=0.2em}

% ══════════════════════════════════════════════════════════════
\begin{document}
\thispagestyle{firstpage}

% ── TITLE BLOCK ───────────────────────────────────────────────
\begin{center}
  \vspace*{0.3cm}

  {\LARGE\bfseries
    The OMARWA--$E_8$ Correspondence:\\[0.3cm]
    A Research Program Connecting\\[0.3cm]
    Modular Periodicity to Exceptional Lie Structure}

  \vspace{0.8cm}

  {\large \"Omer \c{C}etinta\c{s}}\\[0.3cm]
  {\small ORCID: \href{https://orcid.org/0009-0003-2977-6048}{0009-0003-2977-6048}}\\[0.15cm]
  {\small\textit{Independent Researcher}}

  \vspace{0.6cm}
  {\small March 2026}
\end{center}

\vspace{0.6cm}

% ── ABSTRACT ──────────────────────────────────────────────────
\begin{abstract}
We present a research program exploring the correspondence between
the modular periodicity of the integer sequence $T_k = 6 \cdot 2^k + 1$
(OEIS A004119) and the algebraic structure of the exceptional Lie
group~$E_8$. The OMARWA super-period $L = \lcm(P(9), P(11)) = 30$
coincides with the Coxeter number $h(E_8) = 30$, established through
three independent derivations and formally verified in Lean~4
\cite{omarwacore,paket1}. We catalog a series of numerical
correspondences---including Weyl exponent containment
($T_0 = 7, T_1 = 13 \in \exp(E_8)$), the root count identity
$|\Delta(E_8)| = 240 = 8L$, and the ``Exceptional Universality''
whereby all five exceptional Coxeter numbers arise from OMARWA
periods---and frame them as \emph{observations} within a structured
research program. We further examine connections to the P622
crystallographic space group (IT\#177), the palindromic 29-1-29
ribbon architecture satisfying $(29+1) = h(E_8)$, and the
$E_8 \to G_{\mathrm{SM}}$ symmetry breaking chain. Each numerical
match is presented as an observation with explicit epistemic
classification, distinguishing proven arithmetic identities
from model-dependent interpretations. We identify
experimental predictions that would be required for the correspondence
to have physical content and list the free parameters of the framework.

\smallskip\noindent
\textbf{Scope and Disclaimer.} This paper presents a \emph{research
program}, not a completed physical theory. The arithmetic core---modular
periods, super-period $L = 30$, and related identities---is formally
proven~\cite{paket2}. The correspondence with $E_8$ involves numerically
verified but algebraically unproven connections. Physical
interpretations are speculative and classified accordingly.
\end{abstract}

\smallskip
\noindent\textbf{Keywords:} $E_8$ Lie algebra $\cdot$ Coxeter number
$\cdot$ Weyl exponents $\cdot$ OMARWA sequence $\cdot$ modular
periodicity $\cdot$ P622 space group $\cdot$ crystallographic symmetry
$\cdot$ research program

\smallskip
\noindent\textbf{MSC 2020:} 17B25, 20F55, 11B50, 11A07, 20H15

\vspace{0.5cm}

% ══════════════════════════════════════════════════════════════
\section{Introduction and Motivation}
\label{sec:intro}
% ══════════════════════════════════════════════════════════════

The sequence $T_k = 6 \cdot 2^k + 1$ for $k \geq 0$ produces the terms
$7, 13, 25, 49, 97, 193, \ldots$ (OEIS A004119). Despite its elementary
definition, the modular periodicity of this sequence---studied in
detail in~\cite{paket2}---exhibits a remarkably structured theory
culminating in a super-period $L = 30$.

The number~30 is also the Coxeter number of $E_8$, the largest
exceptional simple Lie algebra. This coincidence, together with
a cluster of additional numerical correspondences cataloged below,
motivates the present research program: \emph{is the agreement between
the OMARWA super-period and the $E_8$ Coxeter number accidental, or
does it reflect a deeper algebraic structure?}

We emphasize at the outset that:
\begin{enumerate}
  \item The \textbf{arithmetic core} (modular periods, fractal laws,
    super-period derivation) is \emph{proven}---both by classical methods
    and by formal verification in Lean~4 with Mathlib
    (956~theorems, 28~modules, 0~\texttt{sorry})~\cite{paket1}.

  \item The \textbf{$E_8$ correspondences} cataloged here are
    \emph{observed} numerical coincidences. They are organized
    systematically but are not proven consequences of a unifying
    algebraic mechanism.

  \item The \textbf{physical interpretations}---symmetry breaking scales,
    generation structure, experimental predictions---are \emph{speculative}
    and presented as such.
\end{enumerate}

\medskip
\noindent\textbf{Organization.}
Section~\ref{sec:review} reviews the OMARWA period theory.
Section~\ref{sec:e8} recalls the structure of~$E_8$.
Section~\ref{sec:correspondence} catalogs the numerical correspondences.
Section~\ref{sec:p622} treats the P622 crystallographic connection.
Section~\ref{sec:ribbon} develops the 29-1-29 palindromic ribbon.
Section~\ref{sec:breaking} outlines the $E_8 \to G_{\mathrm{SM}}$ chain.
Section~\ref{sec:gs} presents the Grand Synthesis observations.
Section~\ref{sec:predictions} states experimental predictions.
Section~\ref{sec:free} lists free parameters.
Section~\ref{sec:conclusion} discusses open questions.

% ══════════════════════════════════════════════════════════════
\section{Review of the OMARWA Period Theory}
\label{sec:review}
% ══════════════════════════════════════════════════════════════

We recall the principal results established in~\cite{paket2} and
formally verified in~\cite{paket1}.

\begin{definition}[OMARWA Sequence]
$T_k = 6 \cdot 2^k + 1$ for $k \geq 0$.
The affine recurrence is $T_{k+1} = 2T_k - 1$, $T_0 = 7$.
\end{definition}

\begin{definition}[OMARWA Period]
For $m \geq 1$, the period $P(m)$ is the minimal positive integer~$p$
such that $T_{k+p} \equiv T_k \pmod{m}$ for all $k \geq 0$.
\end{definition}

\begin{theorem}[Core Reduction {\cite[Thm.~3.1]{paket2}}]
\label{thm:core}
$P(m) = \ord_{m'}(2)$ where $m' = m / \gcd(m, 6)$.
\end{theorem}

\begin{theorem}[Super-Period {\cite[Thms.~3.9--3.11]{paket2}}]
\label{thm:super}
The super-period is $L = \lcm(P(9), P(11)) = \lcm(2, 10) = 30$,
derived through three independent paths:
\begin{enumerate}
  \item $L = \lcm(P(9), P(11)) = \lcm(2, 10) = 30$.
  \item $L = \lcm(P(9), P(11), P(27)) = \lcm(2, 10, 6) = 30$.
  \item $L = P(77) = \ord_{77}(2) = 30$ (direct computation).
\end{enumerate}
\end{theorem}

\begin{theorem}[Fractal Period Law {\cite[Thm.~3.5]{paket2}}]
\label{thm:fractal}
For non-Wieferich primes $p \nmid 6$:
$P(p^n) = \ord_p(2) \cdot p^{n-1}$.
In particular, $P(3^n) = 2 \cdot 3^{n-2}$ for $n \geq 2$ and
$P(11^n) = 10 \cdot 11^{n-1}$.
\end{theorem}

\begin{theorem}[Exceptional Universality {\cite[Thm.~11.1]{paket2}}]
\label{thm:exceptional}
All five exceptional Lie algebra Coxeter numbers are recoverable from
the OMARWA periods $P(9)$, $P(11)$, $P(27)$ and the P622 point group
order $|D_6|$:

\smallskip
\begin{center}
\begin{tabular}{@{}cccl@{}}
\toprule
Group & $h$ & OMARWA Origin & Derivation \\
\midrule
$G_2$ & 6  & $P(27) = 6$ & $\ord_9(2) \cdot 3$\\
$F_4$ & 12 & $|D_6| = 12$ & P622 point group order \\
$E_6$ & 12 & $|D_6| = 12$ & Same \\
$E_7$ & 18 & $P(11) + P(9) + P(27) = 18$ & Period sum \\
$E_8$ & 30 & $L = \lcm(P(9), P(11), P(27))$ & Super-period \\
\bottomrule
\end{tabular}
\end{center}
\end{theorem}

\begin{remark}
All theorems in this section carry epistemic class~\textbf{[A]}:
they are formally established and machine-verified. The computational
datasets supporting these results are available as~\cite{paket3}.
\end{remark}

% ══════════════════════════════════════════════════════════════
\section{The $E_8$ Exceptional Lie Algebra}
\label{sec:e8}
% ══════════════════════════════════════════════════════════════

We recall standard facts about $E_8$; all material in this section is
classical~\cite{humphreys,bourbaki}.

\subsection{Classification Context}

The Killing--Cartan classification of simple Lie algebras yields four
infinite families ($A_n$, $B_n$, $C_n$, $D_n$) and five exceptionals
($G_2$, $F_4$, $E_6$, $E_7$, $E_8$). The algebra~$E_8$ is the largest:

\begin{center}
\begin{tabular}{@{}ccccc@{}}
\toprule
Algebra & $\dim$ & $\rank$ & $h$ & $|\Delta|$ \\
\midrule
$G_2$   & 14   & 2 & 6  & 12  \\
$F_4$   & 52   & 4 & 12 & 48  \\
$E_6$   & 78   & 6 & 12 & 72  \\
$E_7$   & 133  & 7 & 18 & 126 \\
$E_8$   & 248  & 8 & 30 & 240 \\
\bottomrule
\end{tabular}
\end{center}

\subsection{Dynkin Diagram and Cartan Matrix}

The $E_8$ Dynkin diagram has 8 nodes with branching at $\alpha_3$:

\begin{center}
$\alpha_1 - \alpha_3 - \alpha_4 - \alpha_5 - \alpha_6 -
\alpha_7 - \alpha_8$, \quad with $\alpha_2$ branching from $\alpha_3$.
\end{center}

\noindent The Cartan matrix $A_{E_8}$ is the $8\times 8$ symmetric
matrix with $\det(A_{E_8}) = 1$ (unimodular), making the root and
weight lattices identical: $P(E_8) = Q(E_8) = \Lambda_{E_8}$.

\subsection{Root System: 240 Vectors in $\mathbb{R}^8$}

The 240 root vectors divide into two classes:
\begin{itemize}
  \item \textbf{Class I} ($112$ roots): all permutations of
    $(\pm 1, \pm 1, 0, 0, 0, 0, 0, 0)$.
  \item \textbf{Class II} ($128$ roots):
    $(\pm\tfrac12,\ldots,\pm\tfrac12)$ with an even
    number of minus signs.
\end{itemize}
Total: $112 + 128 = 240$. These form the densest sphere packing in
$\mathbb{R}^8$~\cite{viazovska}, with kissing number~240.

\subsection{Weyl Group and Exponents}

The Weyl group $W(E_8)$ has order
$|W(E_8)| = 696\,729\,600 = 2^{14}\cdot 3^5 \cdot 5^2 \cdot 7$.
The eight Weyl exponents are:

\begin{equation}
\label{eq:exponents}
\boxed{\exp(E_8) = \{1, 7, 11, 13, 17, 19, 23, 29\}.}
\end{equation}

These are precisely the integers in $\{1,\ldots,29\}$ coprime to~30:
$\exp(E_8) = \{m \in \{1,\ldots,29\} : \gcd(m, 30) = 1\}$.
Since $\varphi(30) = 8 = \rank(E_8)$, this is the \textbf{Euler
totient identity for~$E_8$}.

\subsection{The Coxeter Number $h = 30$}

The highest root of $E_8$ is
$\theta = 2\alpha_1 + 3\alpha_2 + 4\alpha_3 + 6\alpha_4 + 5\alpha_5 +
4\alpha_6 + 3\alpha_7 + 2\alpha_8$, with coefficient sum
$2+3+4+6+5+4+3+2 = 29$, yielding:

\begin{equation}
\label{eq:coxeter}
h(E_8) = 29 + 1 = 30.
\end{equation}

The Coxeter number admits at least five equivalent characterizations:
height of highest root plus one, order of the Coxeter element,
$|\Delta|/\rank = 240/8$, $1 + \max(\exp)$, and the dual Coxeter
number $h^\vee$ (equal to $h$ since $E_8$ is simply laced).

% ══════════════════════════════════════════════════════════════
\section{The Numerical Correspondences}
\label{sec:correspondence}
% ══════════════════════════════════════════════════════════════

We now catalog the observed correspondences between OMARWA arithmetic
and $E_8$ structure. Each is stated as an \textbf{Observation},
explicitly distinguishing it from a proven theorem.

\subsection{The Coxeter Coincidence}

\begin{observation}[Coxeter Identity]
\label{obs:coxeter}
The OMARWA super-period equals the $E_8$ Coxeter number:
\begin{equation}
L = \lcm(P(9), P(11)) = 30 = h(E_8).
\end{equation}
\end{observation}

The left-hand side is a number-theoretic quantity (proven, [A]);
the right-hand side is a Lie-theoretic quantity (classical, [A]).
The equality $30 = 30$ is trivially true. The non-trivial question is
whether there exists an algebraic mechanism connecting them---i.e.,
whether there is a functor or homomorphism from the OMARWA period
lattice to the $E_8$ Coxeter theory.

\textbf{Status:} The arithmetic identity is proven; the structural
connection is an open question.

\subsection{Weyl Exponent Containment}

\begin{observation}[Exponent Overlap]
\label{obs:weyl}
The first two OMARWA terms lie in the Weyl exponent set:
\begin{equation}
T_0 = 7 \in \exp(E_8), \qquad T_1 = 13 \in \exp(E_8).
\end{equation}
\end{observation}

Since the exponents are the integers coprime to~30 in $[1,29]$, and
since $\gcd(7, 30) = 1$ and $\gcd(13, 30) = 1$, this containment is
arithmetically transparent. The third term $T_2 = 25$ fails
($\gcd(25, 30) = 5$), and $T_3 = 49 > 29$ exceeds the range.

\textbf{Status:} Arithmetic fact [A]. The question is whether the
OMARWA seed $T_0 = 7$ and $T_1 = 13$ play a distinguished role
within $E_8$ representation theory.

\subsection{Root Count Identity}

\begin{observation}[Root--Period Correlation]
\label{obs:roots}
The number of $E_8$ roots factorizes through the OMARWA super-period:
\begin{equation}
|\Delta(E_8)| = 240 = 8 \times 30 = \rank(E_8) \times L.
\end{equation}
\end{observation}

This is an instance of the general identity
$|\Delta(G)| = \rank(G) \cdot h(G)$ valid for all simply-laced Lie
algebras~\cite{humphreys}. Given Observation~\ref{obs:coxeter},
the factorization through~$L$ follows. The observation is therefore
a corollary rather than an independent coincidence.

\textbf{Status:} Corollary of Observation~\ref{obs:coxeter} and
classical Lie theory [A].

\subsection{Euler Totient Identity}

\begin{observation}[Totient--Rank]
\label{obs:totient}
$\varphi(L) = \varphi(30) = 8 = \rank(E_8)$.
\end{observation}

Since the exponents of a simply-laced Lie algebra of rank~$r$ are the
$r$ integers coprime to $h$ in $[1, h-1]$, this identity is equivalent
to the statement that $E_8$ is simply laced with $h = 30$ and $r = 8$.

\textbf{Status:} Corollary of classical theory [A].

\subsection{The Sabbath Node and Mersenne Periods}

\begin{observation}[Sabbath--Rank Bridge]
\label{obs:sabbath}
The period of the Mersenne prime $M_5 = 31$ satisfies
$P(31) = 5$. The period of the Fermat prime $F_3 = 17$ satisfies
$P(17) = 8 = \rank(E_8)$.
\end{observation}

The term $T_6 = 385 = 5 \times 7 \times 11$ is the ``Sabbath node''
where three OMARWA primes converge~\cite{paket2}. This triple
factorization yields $P(385) = \lcm(P(5), P(7), P(11)) =
\lcm(4, 3, 10) = 60 = 2L$.

\textbf{Status:} Arithmetic facts [A]. Physical interpretation [C].

\subsection{Dimensional Decomposition}

\begin{observation}[Adjoint Dimension]
\label{obs:dim}
$\dim(\mathfrak{e}_8) = 248 = |\Delta| + \rank = 240 + 8 = 8(L+1) =
8 \times 31$.
\end{observation}

The factor $31 = 2^5 - 1$ is a Mersenne prime, and $P(31) = 5$ in
the OMARWA theory.

\textbf{Status:} Arithmetic identity [A]. Whether the ``$L+1$''
structure has Lie-theoretic content beyond numerology is open.

% ══════════════════════════════════════════════════════════════
\section{The P622 Crystallographic Connection}
\label{sec:p622}
% ══════════════════════════════════════════════════════════════

\subsection{Space Group P622}

The crystallographic space group P622 (Hermann--Mauguin designation)
has International Tables number \textbf{177}~\cite{hahn}. It belongs
to the hexagonal crystal system with point group $D_6$ of order~12.
P622 is the maximal Sohncke group in the hexagonal system---the
largest chiral-compatible space group.

\begin{center}
\begin{tabular}{@{}ll@{}}
\toprule
Property & Value \\
\midrule
Hermann--Mauguin & P622 \\
IT number & 177 \\
Point group & $D_6$ (622) \\
$|D_6|$ & 12 \\
Crystal system & Hexagonal \\
Lattice & Primitive (P) \\
Enantiomorphic children & P$6_122$ (178), P$6_522$ (179) \\
\bottomrule
\end{tabular}
\end{center}

\subsection{The 354--177 Identity}

\begin{observation}[IT Number--Half Period]
\label{obs:p622}
The OMARWA holon count $354 = 6 \times 59$ satisfies:
\begin{equation}
\frac{354}{2} = 177 = \mathrm{IT\#}(\text{P622}).
\end{equation}
\end{observation}

The factorization $354 = 2 \times 3 \times 59$ arises from the
six-sector holon hierarchy ($6 \times 59$ channels). The Core
Reduction Theorem gives $P(354) = \ord_{59}(2) = 58$, and equally
$P(177) = \ord_{59}(2) = 58$, so both carry the same OMARWA period.

\textbf{Status:} Arithmetic [A]; crystallographic identification [B].
The question is whether the appearance of IT\#177 is a meaningful
structural link or a numerical accident.

\subsection{Point Group Order and Gauge Dimension}

\begin{observation}[Dimensional Coincidence]
\label{obs:gauge}
$|D_6| = 12 = \dim(G_{\mathrm{SM}})$ where
$G_{\mathrm{SM}} = SU(3)_C \times SU(2)_L \times U(1)_Y$ with
$\dim = 8 + 3 + 1 = 12$.
\end{observation}

\textbf{Status:} Numerical fact [A]. Physical interpretation~[C].

\subsection{The Enantiomorphic Pair and Chirality}

P622 is the achiral parent of the enantiomorphic pair
P$6_122$ (IT\#178) and P$6_522$ (IT\#179):
\begin{equation}
\text{P622} \supset \{\text{P}6_122,\; \text{P}6_522\}.
\end{equation}

The $6_1$ and $6_5$ screw axes define left- and right-handed
helical structures. In the OMARWA framework, this chirality is
realized through the 29-1-29 palindromic ribbon (Section~\ref{sec:ribbon}),
where the left and right halves correspond to the two enantiomorphic
sub-groups.

% ══════════════════════════════════════════════════════════════
\section{The 29-1-29 Palindromic Ribbon}
\label{sec:ribbon}
% ══════════════════════════════════════════════════════════════

\subsection{Construction}

Each of the six sectors contains 59 channels. The integer 59
admits a palindromic decomposition:

\begin{equation}
59 = 29 + 1 + 29.
\end{equation}

\begin{definition}[Sector Ribbon]
The 59-channel sector sequence is
\[
S_{59} = (L_{29}, L_{28}, \ldots, L_1,\; U_0,\; R_1, \ldots, R_{29}),
\]
where $L_j$ are left-chiral channels, $R_j$ right-chiral channels,
and $U_0$ is the central Unity Axis.
\end{definition}

The reflection operator $\mathcal{R}(x_i) = x_{59-i}$ is an
involution fixing $U_0$, mapping $L_j \leftrightarrow R_j$.

\subsection{Number-Theoretic Properties}

\begin{proposition}[Sophie Germain Pair]
29 is a Sophie Germain prime and $59 = 2 \times 29 + 1$ is its
associated safe prime.
\end{proposition}

\begin{proposition}[Uniqueness]
\label{prop:unique}
The decomposition $59 = 29 + 1 + 29$ is the unique palindromic
partition with simultaneously prime wings and trivial center
satisfying $(29 + 1) = h(E_8)$.
\end{proposition}

\begin{proof}
Any palindrome $59 = a + c + a$ with $c = 1$ forces $a = 29$
(prime). Other decompositions either have composite wings
($28 = 2^2 \times 7$, $27 = 3^3$, etc.) or non-unit center,
and none satisfy $(a + 1) = 30$.
\end{proof}

\subsection{The Coxeter Window}

Grouping adjacent elements $(1 + 29) = 30$ across sector boundaries
yields the \textbf{Coxeter window}:

\begin{observation}[Coxeter Window]
\label{obs:window}
$(29 + 1) = h(E_8) = 30$.
\end{observation}

Since there are 12 such windows around the full circle:
$12 \times 30 = 360 = 354 + 6$.

\subsection{Weyl Exponent as Half-Ribbon Width}

\begin{observation}[Maximal Exponent]
\label{obs:maxexp}
The half-ribbon width $|L| = |R| = 29 = e_8$, the maximal Weyl
exponent of $E_8$.
\end{observation}

The Poincar\'e polynomial of $E_8$ has maximal factor
$(t^{60}-1)/(t^2-1)$ of degree $2 \times 29 = 58$, matching the
total chiral channel count $|L| + |R| = 58$ per sector.

\subsection{Modular Periods of the Ribbon}

From the Core Reduction Theorem:
\begin{itemize}
  \item $P(59) = \ord_{59}(2) = 58 = \varphi(59)$: full ribbon
    traversal (2 is a primitive root mod~59).
  \item $P(29) = \ord_{29}(2) = 28 = \varphi(29)$: half-ribbon
    traversal (2 is a primitive root mod~29).
  \item $P(354) = P(177) = \ord_{59}(2) = 58$: period bridge.
\end{itemize}

\begin{theorem}[Full Ribbon Traversal {\cite[Thm.~9.2]{paket2}}]
The OMARWA sequence visits every chiral channel of the 59-ribbon
exactly once per period:
$|\{T_k \bmod 59 : 0 \leq k < 58\}| = 58 = |S_{59}| - 1$.
The Unity Axis position ($r \equiv 1 \pmod{59}$) is the unique
unvisited residue.
\end{theorem}

% ══════════════════════════════════════════════════════════════
\section{The $E_8 \to G_{\mathrm{SM}}$ Symmetry Breaking Chain}
\label{sec:breaking}
% ══════════════════════════════════════════════════════════════

\subsection{The Canonical Sequence}

The standard $E_8$ breaking chain in grand unification
theory~\cite{slansky,georgi,gursey} proceeds via:

\begin{equation}
\label{eq:chain}
\begin{split}
E_8 &\xrightarrow{\text{SSB}_1}
E_6 \times SU(3) \xrightarrow{\text{SSB}_2}
SO(10) \times U(1) \\
&\xrightarrow{\text{SSB}_3}
SU(5) \times U(1) \xrightarrow{\text{SSB}_4}
G_{\mathrm{SM}} \xrightarrow{\text{SSB}_5}
SU(3)_C \times U(1)_{\mathrm{EM}}.
\end{split}
\end{equation}

\subsection{Representation Decomposition}

At the first step, the adjoint decomposes as:
\begin{equation}
248 = (78, 1) \oplus (27, 3) \oplus (\overline{27}, \overline{3})
\oplus (1, 8).
\end{equation}

The $(27, 3)$ component yields three copies of the $E_6$
fundamental representation---one per fermion generation. The fact
that exactly three generations emerge from the $SU(3)$ factor is
a representation-theoretic consequence of the branching
rule~\cite{gursey,slansky}.

\subsection{OMARWA Scale Assignment}

\begin{observation}[Scale Indexing]
\label{obs:scale}
Each breaking step can be indexed by an OMARWA term:

\smallskip
\begin{center}
\begin{tabular}{@{}ccccc@{}}
\toprule
Step & Breaking & $T_k$ & Scale (GeV) & Class \\
\midrule
SSB$_1$ & $E_8 \to E_6 \times SU(3)$ & $T_0 = 7$ & $\sim 10^{19}$ & [C] \\
SSB$_2$ & $E_6 \to SO(10) \times U(1)$ & $T_1 = 13$ & $\sim 10^{16}$ & [C] \\
SSB$_3$ & $SO(10) \to SU(5) \times U(1)$ & $T_2 = 25$ & $\sim 10^{13}$ & [C] \\
SSB$_4$ & $SU(5) \to G_{\mathrm{SM}}$ & $T_3 = 49$ & $\sim 10^{10}$ & [C] \\
SSB$_5$ & EW breaking & $T_5 = 193$ & $\sim 10^{2}$ & [C] \\
\bottomrule
\end{tabular}
\end{center}
\end{observation}

\textbf{Status:} All entries are epistemic class~[C]---speculative
assignments motivated by the numerical sequence but without
dynamical justification.

\subsection{Three-Generation Structure}

\begin{observation}[Generation--Scale Correspondence]
\label{obs:gen}
Each fermion generation can be associated with an OMARWA scale:

\smallskip
\begin{center}
\begin{tabular}{@{}ccccc@{}}
\toprule
Gen. & Fermions & $T_k$ & Modular Period & Class \\
\midrule
1st & $e, \nu_e, u, d$ & $T_0 = 7$ & $P(9) = 2$ & [C] \\
2nd & $\mu, \nu_\mu, c, s$ & $T_1 = 13$ & $P(11) = 10$ & [C] \\
3rd & $\tau, \nu_\tau, t, b$ & $T_2 = 25$ & $P(27) = 6$ & [C] \\
\bottomrule
\end{tabular}
\end{center}
\end{observation}

\textbf{Status:} Class [C]. The representation-theoretic origin of
three generations (from the $SU(3)$ factor in~\eqref{eq:chain}) is
classical [A]. The specific OMARWA scale assignments are conjectural.

% ══════════════════════════════════════════════════════════════
\section{Grand Synthesis Observations}
\label{sec:gs}
% ══════════════════════════════════════════════════════════════

The following observations, drawn from the OMARWA computational
analysis and formally verified in Lean~4 where indicated
(module: \texttt{GrandSynthesis.lean}~\cite{paket1}), consolidate
the multi-domain numerical correspondences.

\subsection{GS-1: Fibonacci Residues in the mod~11 Orbit}

\begin{observation}[Fibonacci--mod~11 Bridge]
\label{obs:gs1}
The OMARWA residues modulo~11 satisfy
$T_1 \bmod 11 = 2 = F_3$, $T_2 \bmod 11 = 3 = F_4$,
$T_3 \bmod 11 = 5 = F_5$---three consecutive Fibonacci numbers.
\end{observation}

\textbf{Lean~4 status:} Verified (\texttt{gs1\_fibonacci\_orbit}). Class [A].

\subsection{GS-3: CRT Pascal Decomposition and P622}

\begin{observation}[CRT--Crystal Bridge]
\label{obs:gs3}
$T_k \bmod 6 = (T_k \bmod 2, T_k \bmod 3)$ via CRT.
Since $T_k \equiv 1 \pmod{2}$ and $T_k \equiv 1 \pmod{3}$ for all~$k$,
the mod~6 Pascal decomposition yields the fixed point $(1,1)$,
which maps to the P622 crystallographic identity element.
\end{observation}

\textbf{Lean~4 status:} Verified. Class [A] (arithmetic); [B] (crystal interpretation).

\subsection{GS-7: Lucas--Golden Bridge}

\begin{observation}[Lucas--Fibonacci Convergence]
\label{obs:gs7}
$L_5 = 11$, $L_7 = 29$, and $B_2 = 55 = F_{10} = F_5 \times L_5$,
exhibiting a triple convergence of OMARWA-related integers, Fibonacci
numbers, and Lucas numbers.
\end{observation}

Here $B_2 = \sum_{k=0}^{2} T_k - 3 = 7 + 13 + 25 - 3 \cdot 1 = 42$
\ldots{} [Note: $B_2$ uses a specific OMARWA block-sum
convention~\cite{paket2}.]

\textbf{Lean~4 status:} Verified (\texttt{gs7\_lucas\_golden\_bridge}). Class [A].

\subsection{GS-9: Palindromic Wing Sums and $h(E_8)/2$}

\begin{observation}[Apex Sum Identity]
\label{obs:gs9}
In the mod~30 residue cycle, the palindromic digital root wings sum to
$26 + 28 = 54 = S(11)$, while the apex elements sum to
\begin{equation}
9 + 6 = 15 = \frac{h(E_8)}{2}.
\end{equation}
\end{observation}

\textbf{Lean~4 status:} Verified (\texttt{gs9\_palindromic\_wings}). Class [A] (arithmetic);
[B] (E$_8$ interpretation).

\subsection{GS-10: Master Five-Way Bridge}

\begin{observation}[Five-Way Convergence]
\label{obs:gs10}
The following five independently arising structures all interface
through the super-period $L = 30$:
\begin{enumerate}
  \item Modular period lattice: $L = \lcm(P(9), P(11))$.
  \item $E_8$ Coxeter number: $h(E_8) = 30$.
  \item Fibonacci--Lucas: $F_5 \times L_5 = 5 \times 11 = 55$, with
    $\varphi(30) = 8 = \rank(E_8)$.
  \item Palindromic ribbon: $(29 + 1) = 30$.
  \item CRT--P622 crystallography: $|D_6| = 12$, $12 \times 30 = 360$.
\end{enumerate}
\end{observation}

\textbf{Lean~4 status:} Verified (\texttt{gs10\_master\_synthesis}). Class [A] (individual identities); [B] (unified interpretation).

% ══════════════════════════════════════════════════════════════
\section{Experimental Predictions and Falsifiability}
\label{sec:predictions}
% ══════════════════════════════════════════════════════════════

For the OMARWA--$E_8$ correspondence to constitute a physical theory
rather than a numerical pattern, it must generate falsifiable
predictions. We list the principal predictions and the conditions
under which the framework would be refuted.

\subsection{What Would Need to Be True}

For the correspondence to have physical content:
\begin{enumerate}
  \item The $E_8$ breaking chain~\eqref{eq:chain} must be realized
    in nature (not merely in string compactification models).
  \item The OMARWA scale assignments (Observation~\ref{obs:scale})
    must correspond to actual symmetry-breaking thresholds.
  \item The three-generation structure must originate from the
    $SU(3)$ factor in $E_8 \supset E_6 \times SU(3)$, not from
    topology of the compact manifold.
  \item The P622 crystallographic structure must manifest in
    low-energy physics beyond dimensional coincidence.
\end{enumerate}

\subsection{Collider Predictions}

\begin{observation}[BSM Scale]
\label{obs:bsm}
If $M_{\mathrm{BSM}} = v_{\mathrm{EW}} \cdot \sqrt{h(E_8)}
= 246\;\mathrm{GeV} \times \sqrt{30} \approx 1.35\;\mathrm{TeV}$,
then new physics should appear at HL-LHC energies.
\end{observation}

\textbf{Status:} Class [C]. This is a parameter-free prediction: the
only inputs are the electroweak VEV and the Coxeter number. It is
falsifiable by HL-LHC null results above $1.5\;\mathrm{TeV}$.

\subsection{Neutrino Sector}

\begin{observation}[CP Phase]
\label{obs:cp}
The OMARWA--Coxeter structure suggests $\delta_{CP} = 1.36\pi$
for the Dirac CP-violation phase in the PMNS matrix.
\end{observation}

\textbf{Status:} Class [C]. Testable by DUNE and Hyper-Kamiokande.
Current experimental data~\cite{nufit} constrain
$\delta_{CP} \in [0.8\pi, 2.0\pi]$ at $3\sigma$.

\subsection{Precision Tests}

\begin{observation}[Muon $g-2$]
\label{obs:gm2}
The $E_8$ spectral decomposition predicts a BSM contribution
$\Delta a_\mu^{\mathrm{BSM}} \approx 260 \times 10^{-11}$,
within the Fermilab measurement $249 \pm 48 \times 10^{-11}$.
\end{observation}

\textbf{Status:} Class [C]. Consistency with current data is necessary
but not sufficient; Run-4/5 at Fermilab will reduce the uncertainty.

\subsection{Falsification Criteria}

The research program would be significantly weakened if:
\begin{itemize}
  \item HL-LHC observes no BSM physics below $2\;\mathrm{TeV}$.
  \item DUNE measures $\delta_{CP}$ inconsistent with $1.36\pi$ at
    $5\sigma$.
  \item Muon $g-2$ Run-4/5 excludes $\Delta a_\mu$ in the range
    $[200, 320] \times 10^{-11}$.
  \item Proton decay is observed at rates inconsistent with the
    $E_8$ chain lifetime estimates.
\end{itemize}

These criteria are \emph{necessary but not sufficient}: even if all
predictions survive, the correspondence remains a research program
until an algebraic mechanism is identified.

% ══════════════════════════════════════════════════════════════
\section{Free Parameters and Open Structure}
\label{sec:free}
% ══════════════════════════════════════════════════════════════

Transparency about free parameters is essential for any research
program. We distinguish three categories:

\subsection{Zero-Parameter Identities}

The following have no free parameters:
\begin{itemize}
  \item $L = 30 = h(E_8)$ (Observation~\ref{obs:coxeter}).
  \item $T_0 = 7, T_1 = 13 \in \exp(E_8)$ (Observation~\ref{obs:weyl}).
  \item $|\Delta| = 240 = 8L$ (Observation~\ref{obs:roots}).
  \item $(29+1) = 30$ (Observation~\ref{obs:window}).
  \item $354/2 = 177 = \mathrm{IT\#(P622)}$ (Observation~\ref{obs:p622}).
  \item GS-1 through GS-10 (Section~\ref{sec:gs}).
\end{itemize}

\subsection{One-Parameter Predictions}

The BSM scale (Observation~\ref{obs:bsm}) depends only on the
electroweak VEV $v = 246\;\mathrm{GeV}$ and $h(E_8) = 30$---both
fixed---but the functional form $M \propto v\sqrt{h}$ is an
\emph{assumption} (one structural free parameter: the power of~$h$).

\subsection{Model-Dependent Parameters}

The SSB scale assignments (Observation~\ref{obs:scale}) and
generation-scale correspondence (Observation~\ref{obs:gen}) involve
$\geq 5$ free choices (which $T_k$ maps to which breaking step).
These are classified [C] and must be constrained by independent data
before they acquire predictive power.

% ══════════════════════════════════════════════════════════════
\section{Discussion and Open Questions}
\label{sec:conclusion}
% ══════════════════════════════════════════════════════════════

\subsection{Summary of Correspondences}

Table~\ref{tab:summary} collects all observations with their
epistemic classifications.

\begin{table}[H]
\centering
\caption{Summary of OMARWA--$E_8$ correspondences.}\label{tab:summary}
\smallskip
\begin{tabular}{@{}clcc@{}}
\toprule
\# & Observation & Class & Ref. \\
\midrule
1 & $L = 30 = h(E_8)$ & [A]+[A] & Obs.~\ref{obs:coxeter} \\
2 & $T_0, T_1 \in \exp(E_8)$ & [A] & Obs.~\ref{obs:weyl} \\
3 & $|\Delta| = 8L = 240$ & [A] & Obs.~\ref{obs:roots} \\
4 & $\varphi(L) = \rank(E_8) = 8$ & [A] & Obs.~\ref{obs:totient} \\
5 & $P(17) = 8 = \rank(E_8)$ & [A] & Obs.~\ref{obs:sabbath} \\
6 & $354/2 = 177 = \mathrm{IT\#(P622)}$ & [B] & Obs.~\ref{obs:p622} \\
7 & $|D_6| = 12 = \dim(G_{\mathrm{SM}})$ & [C] & Obs.~\ref{obs:gauge} \\
8 & $(29+1) = h(E_8)$ & [A] & Obs.~\ref{obs:window} \\
9 & $e_8 = 29 = $ half-ribbon & [A] & Obs.~\ref{obs:maxexp} \\
10 & SSB scale indexing & [C] & Obs.~\ref{obs:scale} \\
11 & Generation--scale map & [C] & Obs.~\ref{obs:gen} \\
12 & GS-1: Fibonacci mod~11 & [A] & Obs.~\ref{obs:gs1} \\
13 & GS-9: apex $= h/2 = 15$ & [A]/[B] & Obs.~\ref{obs:gs9} \\
14 & GS-10: five-way bridge & [A]/[B] & Obs.~\ref{obs:gs10} \\
15 & $M_{\mathrm{BSM}} \approx 1.35\;\mathrm{TeV}$ & [C] & Obs.~\ref{obs:bsm} \\
16 & $\delta_{CP} = 1.36\pi$ & [C] & Obs.~\ref{obs:cp} \\
17 & $\Delta a_\mu \approx 260 \times 10^{-11}$ & [C] & Obs.~\ref{obs:gm2} \\
\bottomrule
\end{tabular}
\end{table}

\subsection{What Is Proven vs.\ What Is Observed}

The arithmetic identities (Observations 1--5, 8--9, 12--13) are
proven and machine-verified. The crystallographic correspondence
(Observation~6) is numerically exact but lacks an algebraic
derivation. The physical interpretations (Observations 7, 10--11,
15--17) are speculative.

The central open question is:

\begin{quote}
\textit{Does there exist a functor, homomorphism, or categorical
equivalence from the multiplicative structure of
$(\mathbb{Z}/m\mathbb{Z})^*$ (governing OMARWA periods) to the
root system or Weyl group of~$E_8$ that explains the observed
correspondences?}
\end{quote}

If such a mechanism exists, the correspondences would be promoted
from observations to theorems. If not, they remain a remarkable
collection of numerical coincidences.

\subsection{Research Directions}

\begin{enumerate}
  \item \textbf{Algebraic mechanism:} Identify a precise algebraic
    link between $\ord_m(2)$ and Coxeter/Weyl theory.

  \item \textbf{Generalization:} The sequence $T_k^{(a,b)} =
    a \cdot 2^k + b$ generalizes OMARWA. For which $(a,b)$ does
    the super-period equal a Coxeter number of some simple Lie
    algebra? Lean~4 verification (module:
    \texttt{AlphaGeneralization})~\cite{paket1} shows that
    $T_k^{(\alpha)} \equiv 1 \pmod{3}$ for all $\alpha$ and~$k$, but
    the Coxeter matching appears specific to $(a,b) = (6,1)$.

  \item \textbf{Experimental tests:} HL-LHC Run~3/4 data at
    $\sqrt{s} = 14\;\mathrm{TeV}$ can test the BSM scale
    prediction. DUNE first oscillation data (expected 2028--2030)
    can test the $\delta_{CP}$ prediction.

  \item \textbf{Coldea experiment extension:} Ising chain measurements
    on CoNb$_2$O$_6$ near the $E_8$ critical
    point~\cite{coldea} can test logarithmic corrections to mass ratios.

  \item \textbf{Formal verification:} Extend the Lean~4 formalization
    to cover the crystallographic and Lie-algebraic claims. Current
    coverage: 28~modules, 956~theorems. Target: formalize the
    Exceptional Universality theorem in terms of Mathlib's Lie
    algebra API.
\end{enumerate}

% ══════════════════════════════════════════════════════════════
\section*{Acknowledgments}
% ══════════════════════════════════════════════════════════════

The author thanks the Lean community and Mathlib contributors for
the formal verification infrastructure, and CERN/Zenodo for open
access archiving.

% ══════════════════════════════════════════════════════════════
% APPENDICES
% ══════════════════════════════════════════════════════════════

\appendix

\section{Lean~4 Proof Catalog}
\label{app:lean}

The following Lean~4 modules (from
\href{https://github.com/omarwa622/OmarwaCore}{OmarwaCore},
DOI:~\cite{paket1}) verify the claims made in this paper:

\begin{center}
\begin{tabular}{@{}lll@{}}
\toprule
Module & Key Theorems & Section \\
\midrule
\texttt{Sequence} & $T_k$ definition, recurrence & \S\ref{sec:review} \\
\texttt{ModularPeriod} & Core Reduction, $P(m) = \ord_{m'}(2)$ & \S\ref{sec:review} \\
\texttt{SuperPeriod} & $L = 30$, three paths & \S\ref{sec:review} \\
\texttt{FractalLaws} & $P(p^n)$ fractal growth & \S\ref{sec:review} \\
\texttt{WeylExponents} & $T_0, T_1 \in \exp(E_8)$ & \S\ref{sec:correspondence} \\
\texttt{ExceptionalUniversality} & All 5 exceptional $h$ values & \S\ref{sec:correspondence} \\
\texttt{Palindrome} & 29-1-29 involution & \S\ref{sec:ribbon} \\
\texttt{GrandSynthesis} & GS-1 through GS-10 & \S\ref{sec:gs} \\
\texttt{AlphaGeneralization} & $T_k^{(\alpha)} \equiv 1 \pmod{3}$ & \S\ref{sec:conclusion} \\
\bottomrule
\end{tabular}
\end{center}

\noindent Total: 28~modules, 956~theorems, 0~\texttt{sorry}.
Lean version: 4.29.0-rc3. Mathlib4 dependency.

\section{Epistemic Classification System}
\label{app:epistemic}

This paper uses the following classification for all claims:

\begin{center}
\begin{tabular}{@{}cp{10cm}@{}}
\toprule
Class & Meaning \\
\midrule
{[A]} & Mathematically proven (formal or rigorous proof exists). \\
{[B]} & Strongly supported model statement; physically motivated
    but not formally established. \\
{[C]} & Speculative: defined research direction with explicit
    assumptions and falsification criteria. \\
{[D]} & Engineering / application / demonstration level. \\
\bottomrule
\end{tabular}
\end{center}

% ══════════════════════════════════════════════════════════════
% REFERENCES
% ══════════════════════════════════════════════════════════════

\begin{thebibliography}{20}

\bibitem{paket1}
\"O.~\c{C}etinta\c{s},
``OmarwaCore: Lean~4 Formalization of the OMARWA Sequence
$T_k = 6 \cdot 2^k + 1$,''
Zenodo, 2026.
DOI: \href{https://doi.org/10.5281/zenodo.19035570}{10.5281/zenodo.19035570}.

\bibitem{paket2}
\"O.~\c{C}etinta\c{s},
``Modular Period Structure, Fractal Laws, and Algebraic Universality
of the Affine Sequence $T_k = 6 \cdot 2^k + 1$,''
Zenodo (preprint), 2026.
DOI: \href{https://doi.org/10.5281/zenodo.19035608}{10.5281/zenodo.19035608}.

\bibitem{paket3}
\"O.~\c{C}etinta\c{s},
``Computational Verification Data for the OMARWA Sequence:
Period Tables, Digital Root Maps, and GCD Matrices,''
Zenodo (dataset), 2026.
DOI: \href{https://doi.org/10.5281/zenodo.19036063}{10.5281/zenodo.19036063}.

\bibitem{omarwacore}
\"O.~\c{C}etinta\c{s},
``OmarwaCore: Lean~4 Formalization of the OMARWA Sequence,''
GitHub, 2026.
\url{https://github.com/omarwa622/OmarwaCore}.

\bibitem{oeis}
OEIS Foundation,
``A004119: $a(0) = 1$; thereafter $a(n) = 3 \cdot 2^{n-1} + 1$,''
\textit{The On-Line Encyclopedia of Integer Sequences},
\url{https://oeis.org/A004119}.

\bibitem{humphreys}
J.\,E.~Humphreys,
\textit{Reflection Groups and Coxeter Groups},
Cambridge Studies in Advanced Mathematics~29, 1990.

\bibitem{bourbaki}
N.~Bourbaki,
\textit{Lie Groups and Lie Algebras: Chapters 4--6},
Springer, 2002.

\bibitem{viazovska}
M.\,S.~Viazovska,
``The sphere packing problem in dimension~8,''
\textit{Ann.\ Math.}, vol.~185, no.~3, pp.~991--1015, 2017.

\bibitem{hahn}
T.~Hahn (ed.),
\textit{International Tables for Crystallography, Volume~A:
Space-Group Symmetry}, 5th~ed., Springer, 2005.

\bibitem{slansky}
R.~Slansky,
``Group Theory for Unified Model Building,''
\textit{Phys.\ Reports}, vol.~79, pp.~1--128, 1981.

\bibitem{georgi}
H.~Georgi and S.\,L.~Glashow,
``Unity of All Elementary-Particle Forces,''
\textit{Phys.\ Rev.\ Lett.}, vol.~32, pp.~438--441, 1974.

\bibitem{gursey}
F.~G\"ursey, P.~Ramond, and P.~Sikivie,
``A Universal Gauge Theory Model Based on $E_6$,''
\textit{Phys.\ Lett.\ B}, vol.~60, no.~2, pp.~177--180, 1976.

\bibitem{gross}
D.\,J.~Gross, J.\,A.~Harvey, E.~Martinec, and R.~Rohm,
``Heterotic String,''
\textit{Phys.\ Rev.\ Lett.}, vol.~54, pp.~502--505, 1985.

\bibitem{coldea}
R.~Coldea, D.\,A.~Tennant, E.\,M.~Wheeler, \textit{et al.},
``Quantum Criticality in an Ising Chain:
Experimental Evidence for Emergent $E_8$ Symmetry,''
\textit{Science}, vol.~327, pp.~177--180, 2010.

\bibitem{nufit}
I.~Esteban, M.\,C.~Gonzalez-Garcia, M.~Maltoni,
T.~Schwetz, and A.~Zhou,
``The fate of hints: updated global analysis of three-flavor
neutrino oscillations,''
\textit{JHEP}, vol.~09, p.~178, 2020.
\url{http://www.nu-fit.org}.

\bibitem{ireland}
K.~Ireland and M.~Rosen,
\textit{A Classical Introduction to Modern Number Theory}, 2nd~ed.,
Springer GTM~84, 1990.

\bibitem{nakahara}
M.~Nakahara,
\textit{Geometry, Topology and Physics}, 2nd~ed.,
IOP Publishing, 2003.

\end{thebibliography}

\bigskip
\noindent\rule{\textwidth}{0.4pt}

\smallskip
{\small\noindent\textit{Manuscript prepared: March 2026.
This is a research program preprint, not a completed theory.
Arithmetic core: formally verified in Lean~4.29.0-rc3 + Mathlib4
(956~theorems, 28~modules, 0~sorry).}\\
\textit{Foundation:
DOI \href{https://doi.org/10.5281/zenodo.19035570}{10.5281/zenodo.19035570} (Software),
DOI \href{https://doi.org/10.5281/zenodo.19035608}{10.5281/zenodo.19035608} (Preprint),
DOI \href{https://doi.org/10.5281/zenodo.19036063}{10.5281/zenodo.19036063} (Dataset).}\\
\textit{Source:
\href{https://github.com/omarwa622/OmarwaCore}{github.com/omarwa622/OmarwaCore}.}}

\end{document}
